Con \(\theta=T-20\), el balance es
\[
C\theta'=P-H\theta.
\]
Los datos dan \(C=40\,\mathrm{J/K}\), \(H=2\,\mathrm{W/K}\) y \(\tau=C/H=20\,\mathrm s\). Para dos arranques \((P_i,\theta_i)\), con datos compatibles con el modelo, se identifican ambos coeficientes por separado exactamente cuando
\[
P_1\theta_2-P_2\theta_1\ne0.
\]
Si ambos comienzan en el ambiente, las pendientes no informan sobre \(H\); permiten identificar \(C\) si al menos una potencia es no nula. Si ambas potencias también son cero, no identifican ninguno.
